%\tableofcontents
\vspace{1em}
\hrule
\vspace{1em}

\begin{comment}



\section{Some Basics}

The tasks are synthetic retrieval (RULER): find a random number planted somewhere in 8K–128K tokens of text, with distractor numbers nearby. There are 20 prompts per cell.
\begin{itemize}
    \item Synthetic Retrieval: This means the test is artificially constructed rather than being a natural real-world task (like summarizing a real book). A piece of data is intentionally hidden inside a massive wall of text just to see if the AI can find it. 
    \item RULER: This is the specific name of the benchmark framework (developed by researchers to evaluate long-context LLMs). It stands for Retrieval, Understanding, and Long-context Evaluation Robustness.
    \item 8K-128K tokens: This indicates the length of the text context that the model has to process. A token is roughly a word or part of a word, so 8K tokens is about 8,000 words, and 128K tokens is about 128,000 words. The model has to search through this large amount of text to find the hidden number.
    \item Distractor numbers: These are other numbers placed in the text that are not the target number. They are meant to confuse the model and make the task harder. The model has to correctly identify the planted number and ignore the distractors.
    \item 20 prompts per cell: The Matrix (Cells): When testing LLMs, researchers usually create a grid or matrix. For example, the horizontal axis might be the context length (8K, 16K, 32K, 64K, 128K) and the vertical axis might be the depth of where the needle is hidden (0\% depth = top of the document, 50\% = middle, 100\% = bottom). Each intersection on this grid is a "cell." This means that for each configuration or "cell" of the experiment (which could be a specific model, context length, or other parameter), there are 20 different test cases (prompts) that the model is evaluated on. Each prompt is a separate instance of the retrieval task. 
\end{itemize}

Some metrics
\begin{itemize}
    \item valid mean task score
\end{itemize}

10 promots -> 100prompts or 1000 prompts? 
not  head-output error but task score.  The task score is the final evaluation metric that measures how well the model performed on the retrieval task. It is calculated based on whether the model correctly identified the planted number in each prompt. The valid mean task score is the average score across all valid prompts (those that were not filtered out due to errors or other issues). This metric gives an overall sense of the model's performance on the retrieval task across different configurations and contexts.

\end{comment}

\section{Scope of the common rate-allocation view}
\label{app:scope}

This appendix expands the derivation and measurement contract used in the main
paper.  The common action space is deliberately narrow: it allocates
\emph{key-token} rates within an attention head, leaves values exact, and measures
attention-output error.  Prior work already combines quantization and eviction
in deployable systems.  Our claim is instead that a common local surrogate makes
it possible to ask a different empirical question: when does the discrete
interior improve over uniform quantization and eviction under matched
information and key-bit budgets?

Throughout, \emph{zero rate} means a zero-bit action in the separable planning
surrogate.  It does not mean that set eviction is exactly identical to
independent additive logit noise.  Section~\ref{app:derivation} makes this
distinction explicit.

\section{Output distortion and the zero-rate approximation}
\label{app:derivation}

\subsection{Attention and finite-bit key perturbations}

For a query $q\in\mathbb R^d$, keys $k_i\in\mathbb R^d$, and values
$v_i\in\mathbb R^{d_v}$, define
\begin{align}
 s_i &= q^\top k_i/\sqrt d, &
 a_i &= \frac{\exp(s_i)}{\sum_j\exp(s_j)}, &
 o &= \sum_i a_i v_i .
\end{align}
Quantizing key $i$ changes its logit by $\delta_i$.  Because
\[
 \frac{\partial o}{\partial s_i}=a_i(v_i-o),
\]
the first-order output perturbation is
\begin{equation}
 \Delta o \approx \sum_i a_i(v_i-o)\delta_i .
 \label{eq:app-linear}
\end{equation}
Consequently,
\begin{align}
 \mathbb E\|\Delta o\|^2
 &\approx
 \sum_i a_i^2\|v_i-o\|^2\operatorname{Var}(\delta_i)
 \nonumber\\
 &\quad+
 \sum_{i\ne j}a_i a_j
 (v_i-o)^\top(v_j-o)\operatorname{Cov}(\delta_i,\delta_j).
 \label{eq:app-covariance}
\end{align}
Equation~\eqref{eq:surrogate} follows only when the residual perturbations are
zero mean and the cross-token covariance contribution vanishes.  Softmax
invariance removes a \emph{common} shift,
$\operatorname{softmax}(s+c\mathbf 1)=\operatorname{softmax}(s)$; it does not
remove arbitrary token-dependent bias.

\subsection{Exact set eviction}

Let $E$ be an evicted token set and $A_E=\sum_{i\in E}a_i$.  Renormalizing the
remaining attention weights yields
\begin{align}
 o_{\setminus E}
 &=\frac{o-\sum_{i\in E}a_i v_i}{1-A_E},\\
 o_{\setminus E}-o
 &=-\frac{\sum_{i\in E}a_i(v_i-o)}{1-A_E},\\
 \|o_{\setminus E}-o\|^2
 &=\frac{\left\|\sum_{i\in E}a_i(v_i-o)\right\|^2}
 {(1-A_E)^2}.
 \label{eq:app-exact-eviction}
\end{align}
The denominator and cross-token inner products prevent an exact token-separable
representation.  For one evicted token $j$, the exact error magnitude is
\[
 \frac{a_j}{1-a_j}\|v_j-o\|,
\]
which is the CAOTE singleton score~\citep{goel2025caote}.  ReST-KV further models
attention redistribution and output reconstruction~\citep{an2026restkv}; our
local planner does not capture those effects.

\subsection{The planning approximation}

When the evicted mass is small and cross terms are neglected,
Eq.~\eqref{eq:app-exact-eviction} motivates
\[
 \widetilde D_0(E)=\sum_{i\in E} w_i,
 \qquad
 w_i=a_i^2\|v_i-o\|^2.
\]
We choose units in which the absolute zero-rate distortion is $d_0=1$ and write
the finite-rate costs as $d_b=\operatorname{Var}(\delta_b)$ in the same logit
units.  Thus
\[
 \widetilde D(\mathbf b)=\sum_i w_i d_{b_i},
 \qquad b_i\in\mathcal B,\quad 0\in\mathcal B.
\]
This convention keeps relative quantizer factors separate from the absolute zero-rate
cost.  It is a local approximation used to choose an allocation; every
reported response is obtained by recomputing attention with the actual
quantized/evicted cache.

\section{Discrete rate allocation and regime coordinates}
\label{app:allocation}

\subsection{The implemented optimizer}

At average key budget $B$, the discrete problem is
\begin{equation}
 \min_{b_i\in\mathcal B}\sum_i w_i d_{b_i}
 \quad\text{subject to}\quad
 \sum_i b_i\le BL.
 \label{eq:app-discrete}
\end{equation}
Introducing multiplier $\lambda\ge0$ separates the token decisions:
\begin{equation}
 b_i^\star(\lambda)
 =\arg\min_{b\in\mathcal B}\{w_i d_b+\lambda b\}.
 \label{eq:app-lagrange}
\end{equation}
We search $\lambda$ to satisfy the budget and use the discrete tiers in
Eq.~\eqref{eq:app-lagrange}; the continuous expression is explanatory, not the
experimental optimizer.

For the relaxation $d_b=\kappa 4^{-b}$,
\[
 b_i^\star
 =\frac12\log_2w_i+c
 =\log_2a_i+\log_2\|v_i-o\|+c.
\]
Adding the value term $\log_2\|v_i-o\|$ changed an \emph{aggregate
ladder-width statistic} by at most $0.035$ bits in our measurements.  This is
not a per-token bound; the end-task allocator therefore uses the full weight
$w_i=a_i^2\|v_i-o\|^2$.

\subsection{Tier extinction}

If $d_b\ge d_0$, tier $b$ consumes positive bits without improving the local
surrogate relative to zero rate and is dominated.  This is a sufficient
extinction condition.  The condition is not necessary: a tier with $d_b<d_0$
can still be dominated if it lies above the lower convex envelope of the other
rate--distortion points.  For the 2-bit tier we report
\[
 D_2=\frac1H\sum_{h=1}^{H}\mathbb I[d_{h,2}\ge d_0].
\]
$D_2$ is measured at a particular model and context.  Although $L$ is absent
from its formula, the head cost distribution changes with $L$, so $D_2$ is not
sequence-length invariant and a new deployment context requires calibration.

\subsection{Why the regime map needs two coordinates}

The sharp side and diffuse side arise for different reasons:
\begin{enumerate}
 \item On the sharp side, large absolute finite-bit noise removes low-bit tiers;
 $D_2$ directly probes this boundary.
 \item On the diffuse side, unequal allocation has little leverage when token
 importances are nearly equal; an importance-dispersion statistic is needed.
\end{enumerate}
For attention-only importance,
$\operatorname{std}_i(\log_2 a_i)=\tau/\ln2$ with $\tau=\operatorname{std}_i(s_i)$,
because softmax subtracts the same log normalizer from every logit.  This
identity alone neither proves that $\tau$ is a complete regime variable nor
establishes a uniform-quantization optimum.  The paper therefore treats the
mixed-to-eviction side as the well-supported boundary;
Table~\ref{tab:app-predictors} compares $D_2$ with dispersion predictors.

\section{Measurement protocol and accounting}
\label{app:protocol}

\subsection{Grid and independent units}

The main (symmetric) grid contains 16 model--context cells from six models and
29,312 query-head--configuration aggregates.  The boundary extension contains
25 cells spanning 4k--256k and 44,416 such aggregates.  Across the six models,
10,240 query-head positions and 2,016 KV-head positions are counted once;
contexts reuse those positions.  Model/cell-level analyses, rather than raw
heads, are therefore the independent statistical units.

PG-19 passages supply natural context.  Each cell is admitted only after the
uncompressed model passes the needle-retrieval validity gate.  The allocation
planner uses its local surrogate, while every reported error comes from actual
key quantization or eviction followed by exact attention-output recomputation.

\subsection{Matched information and budget}

The primary comparison uses $B=3$ bits per key element.  Uniform
quantization stores every key at 3 bits.  The eviction endpoint retains the
number of 8-bit keys allowed by the same key-bit budget.  In the symmetric
contract, the mixed allocator and accumulated-attention eviction both use lagged
importance information; an oracle/current-query result is labeled only as an
upper bound.  The five-corner robustness result changes the reference set and is
reported separately.

This accounting omits value-cache bytes, quantization scales, indices, metadata,
protected windows, alignment, and packing.  Moreover, \emph{8-bit} is the maximum
experimental key tier, not full floating-point precision.  The experiments
therefore establish matched key-bit output error, not matched total memory or
throughput.

\subsection{Responses}

For each head,
\[
 g_h=
 \frac{\min(D_h^{\rm uniform},D_h^{\rm evict})}
 {D_h^{\rm mixed}},
 \qquad
 G_+=\exp\!\left(\frac1H\sum_h\log\max(g_h,1)\right).
\]
$G_+$ is a hindsight portfolio statistic: it clips losses after observing all
three errors.  It is not the performance of a deployed router.  The primary
cell-level response is continuous $G_+$; the fraction with $g_h\ge2$ is a
secondary thresholded band statistic.  Every band result names its corner
and information contracts.

\subsection{Compute resources and software}
\label{app:compute}

All experiments ran as Slurm jobs on one shared academic GPU cluster. 
Every job recorded its GPU and library versions, and the cluster has
NVIDIA H100 80GB HBM3 GPUs, PyTorch 2.13.0, and Hugging Face Transformers
5.16.1 (Python 3.11).  Models were loaded in bfloat16 from pinned local
checkpoints with the hub offline. 

\begin{table}[h]
\centering
\small
\setlength{\tabcolsep}{4pt}
\caption{Hardware per experiment.  Each job ran on one node; wall times are
measured from job records where available.}
\label{tab:app-compute}
\begin{tabular}{@{}p{0.34\linewidth}p{0.25\linewidth}p{0.33\linewidth}@{}}
\toprule
\textbf{Experiment} & \textbf{GPUs per job} & \textbf{Wall time per job} \\
\midrule
Regime map, 7--8B and MoE models (\S\ref{sec:regimes}) & 1$\times$H100 & minutes to a few hours per cell \\
Regime map, Llama-3.3-70B and Qwen3-30B-A3B & 4$\times$H100 (70B), up to 4$\times$H100 (30B) & up to several hours per cell \\
End-task audit, one model--context cell of 20 prompts (\S\ref{sec:endtask}) & 1$\times$H100 & 0.5\,h (8k) to 2.7\,h (128k) \\
Router calibration, 10 prompts & 1$\times$H100 & about 10 minutes \\
Nested-code overhead (App.~\ref{app:endtask}) & 1$\times$H100 & 1.8--2.8\,h per cell \\
Readers, statistics, figures, tests & CPU only & minutes \\
\bottomrule
\end{tabular}
\end{table}

Because the model forward pass is not bit-reproducible across processes even on
one GPU, every comparison that pairs methods prompt by prompt computes all
methods inside one process on the same captured queries, keys, and values;
comparisons across processes are made on cell means only.  The spread is
measurable: the full grid and the paired grid evaluate the same 8k cells in
different processes, with the same prompt indices, needle depths, and question
templates, but with haystack text and needle values drawn from two copies of
the PG-19 corpus that differ.
Their cell means differ by a median of 0.04 and a 90th percentile of 0.15 over
112 cell--method pairs (TurboQuant: at most 0.05); this spread includes both
process nondeterminism and the haystack change.  Per-cell wins and losses
(Table~\ref{tab:app-h2h}) therefore use a 0.15 margin.  Over all 36 cells, the
eight methods shared by the two grids differ in mean score by at most 0.03
(OBCache-K + Ada-KV; \sieve{} 0.002, TurboQuant 0.000): enough to reorder close
eviction methods, but not \sieve{} against eviction or dense quantization.
All end-task numbers in the main text come from the paired grid unless marked
\emph{full grid}.

\section{Additional robustness and deployment evidence}
\label{app:robustness}

\begin{table}[h]
\centering
\small
\setlength{\tabcolsep}{4pt}
\caption{Measured costs of acting on the output-distortion map. Output error only.}
\label{tab:deployment}
\begin{tabular}{@{}p{0.29\linewidth}p{0.25\linewidth}p{0.38\linewidth}@{}}
\toprule
\textbf{Constraint} & \textbf{Measurement} & \textbf{Implication} \\
\midrule
Family label over 4,096 tokens & p90 regret $1.00$ in 6/6 runs (5 cells) & calibrate the route once per context \\
Frozen token assignment & $2.52$--$4.30\times$ lag-1 error & re-budget mixed heads every few steps \\
Shared GQA allocation & $1.000/1.136/1.316\times$ at $n_{\rm rep}=1/4/8$ & route per KV head at a small cost \\
4-bit current-query cascade & closes $82$--$91\%$ (median $89\%$) of lag penalty & re-score from base-tier keys \\
\bottomrule
\end{tabular}
\end{table}


The completed measurements support the following bounded conclusions:
\begin{itemize}
 \item In all 16 main-grid cells, the symmetric lagged-versus-lagged band exceeds the
 all-oracle band by 1.2--12.5 points, but it is lower than the asymmetric
 current-query-versus-lagged numbers.  Five cells are STOP and five NARROW.
 \item Across the 25 extension cells, $D_2$ has Spearman correlation $-0.978$ with both
 $G_+$ and the band; controlling for model, it is $-0.94$ with $G_+$ and
 $-0.915$ with the band.  Similar $D_2$
 values can nevertheless differ by eight band points across architectures.
 \item Prompt-block standard deviation is 0.4--3.0 band points.  Replacing the
 two headline corners by the five-corner minimum lowers the band by 5.1--8.6
 points and can change a categorical verdict.
 \item In six decode-stability runs (five model--context cells), the calibration-time family choice has p90 regret 1.00
 through 4,096 generated tokens, whereas a frozen token allocation costs
 $2.52$--$4.30\times$ a one-step-old allocation.
 \item Enforcing one allocation per shared KV head gives marginal in-band costs
 $1.000$, $1.136$, and $1.316$ at GQA repetition factors 1, 4, and 8.  A 4-bit
 current-query cascade closes 82--91\% (median 89\%) of the lagged-score penalty.
\end{itemize}

\begin{table}[h]
\centering
\small
\caption{Predictor comparison on the 25 regime-map cells (\S\ref{sec:regimes}).  Columns: pooled Spearman correlation with $\gplus$; the same partialled on model identity; leave-one-model-out error of a log-linear fit of $\gplus$, as a multiplicative factor; mean Spearman within the held-out model.  All predictors are cell medians of per-head medians on the rows that define $D_2$; the recomputed $D_2$ matches the reported one to within 0.52 points in every cell.  The heterogeneity statistic is the spread of $\tfrac12\log_2 w_i$; for log-normal $w$ it ranks cells as the AM/GM ratio does.}
\label{tab:app-predictors}
% generated from h0_measurement/bugs/12_paper_main_table/predictor_comparison.csv
\begin{tabular}{@{}lrrrr@{}}
\toprule
predictor & $\rho$ with $\gplus$ & $\rho$ \textbar{} model & held-out model error & within-model $\rho$ \\
\midrule
$D_2$ (ours) & $-0.978$ & $-0.942$ & 1.119$\times$ & $-0.986$ \\
spread of log sensitivity (heterogeneity) & $-0.722$ & $-0.923$ & 1.155$\times$ & $-0.986$ \\
attention concentration $\tau$ & $-0.724$ & $-0.922$ & 1.154$\times$ & $-0.986$ \\
attention entropy & $-0.288$ & $-0.903$ & 1.273$\times$ & $-0.943$ \\
$n_{95}/L$ & $+0.331$ & $+0.021$ & 1.271$\times$ & $+0.286$ \\
context length only & $-0.478$ & $-0.872$ & 1.243$\times$ & $-0.986$ \\
\bottomrule
\end{tabular}

\end{table}

\paragraph{Methods note: fitting the crossing.}
On the 16-cell main grid, the band crosses
35\% between the cells at $D_2=25.8\%$ and $27.0\%$, and the isotonic crossing is
$26.8\%$.  A logistic fit instead places GO at $34.6\%$, outside that bracket,
because the flat tail above 30\% pulls the curve.  We therefore report isotonic
(monotone) crossings, which cannot leave the bracket of measured cells, and use
the logistic only as a flagged cross-check.  On the 25-cell extension the GO
crossing moves to $29.0\%$: Llama-3.3-70B at 96k lies above the GO line at
$D_2=28.8\%$, so models now overlap near the line instead of leaving a gap.

\paragraph{Why the audit uses synthetic retrieval.}
An audit can only rank compression policies on a task where their answers
differ.  We screened a natural task, LongBench-v2~\citep{bai2025longbenchv2} in
forced-choice form, on two development blocks of 52 questions.  Even choosing
the best policy separately for each question beat the best single fixed policy
on only 3 and 2 of the 52 questions, so at these budgets the task barely
separates policies, and its pre-registered gate stopped both blocks before the
confirmation block.  On the RULER tasks of Appendix~\ref{app:endtask}, method
means differ by more than 0.6, so the audit can discriminate there.


\begin{comment}
    
\section{Claim boundary}
\label{app:claims}

\begin{table}[h]
\centering
\small
\begin{tabular}{@{}p{0.29\textwidth}p{0.64\textwidth}@{}}
\toprule
\textbf{Supported} & \textbf{Not yet supported} \\
\midrule
Zero rate is useful inside a stated separable local surrogate
& Exact equivalence between independent quantization noise and set eviction \\
An absolute tier cost can make a finite-bit tier dominated
& A universal one-dimensional phase transition or threshold \\
$D_2$ strongly orders the measured mixed-to-eviction boundary
& A context-free predictor that transfers without calibration \\
The output-error-optimal family label is temporally stable in six measured cells while token allocations become stale quickly
& A robust end-task router, total-KV memory savings, or realized throughput \\
The output-error-optimal allocator is the strongest token-selective policy on the measured retrieval cells
& That it beats dense low-bit quantization in general (it loses on Llama-3.1-8B) \\
Output distortion orders budgets and contexts by difficulty
& That output distortion can choose between method families at a fixed budget \\
Token-granular allocation truncates multi-token answers
& That positional smoothing repairs it (pre-registered test pending) \\
\bottomrule
\end{tabular}
\caption{The empirical and methodological claim boundary of the current draft.}
\label{tab:app-claims}
\end{table}


\end{comment}

\section{End-task audit: protocol, baselines, and full results}
\label{app:endtask}

\paragraph{Protocol.}
Each prompt's context is prefilled at full precision and scored; the cache is
then compressed once, and the question is prefilled \emph{through} the
compressed cache before greedy decoding.  This is the question-agnostic setting
of a prefix cache or multi-turn session.  When the question is inside the
observation window instead, SnapKV scores 1.00 at every budget, because the
window vote points at the needle.  Evaluation uses prompts 100--119; the
\sieve{} router is calibrated on the disjoint prompts 0--9.  A block is valid
when the uncompressed model scores at least 0.9; this excludes Qwen3-8B
multi-value retrieval (0.24 at 8k, 0.04 at 32k).  Scores are RULER string
matches; multi-value and variable tracking give fractional credit.

\paragraph{Budget accounting.}
Every method spends $B$ key code bits per context-key element, verified at run
time from the applied bit maps.  Side information differs by method and is
reported in Table~\ref{tab:app-side} rather than folded into $B$.  Eviction and
\sieve{} keep tokens at 8 bits with the same TurboQuant quantizer, so their
side information scales with the kept fraction.

\paragraph{Storage costs outside the code budget.}
Two further costs are small relative to the differences reported in the main
text.  \emph{Nested tiers.}  A mixed allocator must store tokens at several
widths.  We measured the rate cost of storing the tier ladder
$\{0,3,4,6,8\}$ as one successively refinable 3+1+2+2 code, in which an 8-bit
index truncates to 6, 4, or 3 bits without re-encoding, against separate
Lloyd--Max codebooks, at matched attention-output error; successive refinability is classical~\citep{equitz1991successive}, but our scalar codes are not the asymptotic optimum, so we measured it.  The code itself costs
at most $0.42\%$ rate at any width.  End to end at $B=3$, including the extra
bit that 3-bit tokens keep for the 4-bit cascade score, the overhead is
$+0.34\%$ averaged over four model--context cells (90\% CI $[-0.33,+0.99]$) and
$+2.9\%$ in the worst cell.  A pre-registered extension to four further
cells (Qwen3-30B-A3B and Qwen1.5-MoE-A2.7B at 8k, Mistral-7B at 8k and 32k;
48 prompts), with both codes built in one process from the same captured
tensors, passed every validity check.  At $B=3$ the mean overhead is $+1.5\%$
over the four cells (cell means $-3.6\%$ to $+3.8\%$); at $B=2$ the cell means
are $-1.9\%$ to $+4.9\%$.  No prompt in any cell exceeds the pre-registered
20\% tail at either budget, and the worst single prompt is $+10.2\%$
(Qwen3-30B-A3B, $B=3$).  One secondary, pre-registered mechanism prediction
failed.  We expected no single KV group to carry more than 10\% of a prompt's
error in the two Mistral cells at $B=3$; one Mistral-7B 8k prompt reached
11.6\%, although its overhead stayed small.  Negative values mean the nested code needs fewer bits than
separate codebooks at the same output error.  \emph{Protected window.}  The last 32 context
tokens stay uncompressed in every method: $0.39\%$ of the context at 8k,
$0.10\%$ at 32k, and $0.02\%$ at 128k.  Because it is identical across methods,
it shifts every method's memory equally and does not affect the comparison.

\begin{table}[h]
\centering
\small
\setlength{\tabcolsep}{4pt}
\caption{Baselines as implemented, and side information per key element ($d=128$).}
\label{tab:app-side}
\begin{tabular}{@{}p{0.2\linewidth}p{0.52\linewidth}r@{}}
\toprule
method & what is stored, and how it is chosen & side bits \\
\midrule
TurboQuant-MSE & every token at $B$ bits; random rotation, Lloyd--Max, per-token norm & 0.125 \\
KIVI ($g=32$ / $128$) & every token at $B$ bits; per-channel asymmetric integer, min/max per $g$ tokens, post-RoPE & 1.0 / 0.25 \\
KVQuant & every token at $B$ bits; pre-RoPE per-channel non-uniform codebook; 1\% outliers and token 0 exact & $\approx0.32$ \\
H2O & $\lfloor BL/8\rfloor$ tokens at 8 bits; raw attention summed over all prefill queries, no pooling; no separate recent-token budget beyond the 32-token window shared by all methods & $\approx0.05$ \\
SnapKV & same keep count; 32-query observation window, max-pool 7 & $\approx0.05$ \\
DropKV, Ada-KV, OBCache-K, LaProx & same keep count; published scores; Ada-KV head allocation for Ada-KV and OBCache-K; LaProx global top-$K$ & $\approx0.05$ \\
\sieve{} & per-token width in $\{0,1,2,3,4,5,6,8\}$; Eq.~\ref{eq:discrete} per KV head, 4-bit cascade score; offline per-KV-head router & 0.07--0.09 \\
\bottomrule
\end{tabular}
\end{table}

\paragraph{Deviations from the original papers.}
KIVI keeps its last partial group in full precision; we quantize it as a short
group, and the 32-token uncompressed window shared by all methods plays the
role of KIVI's residual.  
KVQuant calibrates outlier thresholds and its
non-uniform codebook offline with Fisher-weighted $k$-means; we fit both online
on each prompt's own context without weighting.  Fitting on the test context
itself likely favors KVQuant.  The Qwen conclusion does not rest on it: KIVI-128,
which needs no calibration, also beats \sieve{} there on its own (0.986 vs.\ 0.962).  H2O
uses its raw accumulated score, whose causal position bias (a token is seen by
every later query) is H2O's own.  
%All methods use simulated quantization: the model reads dequantized keys, and no packed kernel is involved.

\paragraph{The harder cell (failed gate).}
Every dense quantizer is near its ceiling in the main grid, so we pre-registered
one harder cell: multi-key retrieval with 32 keys, 4 values, and 4 hops at 32k on
Llama-3.1-8B, evaluated on 60 fresh prompts (1000--1059) with a router calibrated
on disjoint prompts.  The cell counts only if the uncompressed model scores at
least 0.95 and TurboQuant at $B=2$ lies in $[0.50,0.90]$.  The uncompressed model
scores 1.00, but TurboQuant scores 0.917, so the gate fails and the cell is not
used as evidence of a non-ceiling regime.  For completeness, at $B=2$ / $B=3$:
KIVI-32 0.98 / 1.00, KIVI-128 0.87 / 0.98, KVQuant 0.82 / 1.00, TurboQuant
0.92 / 1.00; \sieve{} (router) 0.38 / 1.00, interior alone 0.13 / 0.42; best
eviction method H2O 0.20 / 0.42, the others at most 0.15 / 0.38.  With many keys
competing for the kept tokens, eviction collapses while every quantizer stays
near ceiling; \sieve{} lies between them at $B=2$ and matches the quantizers at
$B=3$.

\begin{table}[h]
\centering
\small
\caption{\sieve{} (router) against each baseline over the 36 valid cells.  A win or loss is a cell-mean difference above 0.15, the 90th percentile of the run-to-run spread (App.~\ref{app:compute}); the interval resamples prompts within each model and context.}
\label{tab:app-h2h}
% generated by figures/make_endtask_tables.py from 9 parquets
\begin{tabular}{@{}lrrrrl@{}}
\toprule
\sieve{} (router) vs. & wins & ties & losses & mean $\Delta$ & 90\% CI \\
\midrule
TurboQuant-MSE & 3 & 19 & 14 & $-0.182$ & $[-0.209, -0.154]$ \\
KIVI (g=128) & 0 & 22 & 14 & $-0.202$ & $[-0.231, -0.173]$ \\
KIVI (g=32) & 0 & 22 & 14 & $-0.225$ & $[-0.252, -0.198]$ \\
KVQuant-style & 0 & 22 & 14 & $-0.209$ & $[-0.234, -0.183]$ \\
H2O & 29 & 4 & 3 & $+0.427$ & $[+0.397, +0.459]$ \\
SnapKV & 25 & 9 & 2 & $+0.377$ & $[+0.341, +0.411]$ \\
DropKV & 26 & 9 & 1 & $+0.367$ & $[+0.328, +0.405]$ \\
Ada-KV & 20 & 9 & 7 & $+0.255$ & $[+0.218, +0.290]$ \\
LaProx & 20 & 8 & 8 & $+0.169$ & $[+0.133, +0.204]$ \\
OBCache-K + Ada-KV & 18 & 10 & 8 & $+0.192$ & $[+0.152, +0.230]$ \\
best eviction method per cell & 16 & 10 & 10 & $+0.103$ & --- \\
\bottomrule
\end{tabular}

\end{table}

\begin{table}[h]
\centering
\scriptsize
\setlength{\tabcolsep}{3pt}
\caption{Failure case: Llama-3.1-8B at 128k.  Error types are over single- and multi-key answers at $B\in\{2,3\}$ (80 per method).  \emph{Truncated}: the prediction keeps the correct first three digits (Llama's first number token) but is not the correct number; proper prefixes alone account for 22\% (interior and router), 39\% (H2O), and at most 4\% (pooled eviction methods).  Diagnostic rows are not tuned methods.  \emph{Wrong needle}: another needle's value.}
\label{tab:app-fail}
% generated by figures/make_endtask_tables.py from 9 parquets
\begin{tabular}{@{}lrrrrrrrr@{}}
\toprule
 & \multicolumn{2}{c}{task score} & \multicolumn{3}{c}{single/multi-key answers, $B{=}2,3$} & evicted & head-output & heads to \\
\cmidrule(lr){2-3}\cmidrule(lr){4-6}
method & $B{=}2$ & $B{=}3$ & truncated & wrong needle & other & ($B{=}2$) & error ($B{=}2$) & interior \\
\midrule
TurboQuant-MSE & 0.931 & 0.984 & 0\% & 1\% & 0\% & 0\% & 0.705 & --- \\
KIVI (g=128) & 0.899 & 0.984 & 1\% & 2\% & 0\% & 0\% & 0.447 & --- \\
KIVI (g=32) & 0.949 & 0.997 & 0\% & 0\% & 0\% & 0\% & 0.297 & --- \\
KVQuant-style & 0.904 & 0.971 & 0\% & 0\% & 1\% & 0\% & 0.240 & --- \\
H2O & 0.239 & 0.364 & 52\% & 6\% & 29\% & 75\% & 0.337 & --- \\
SnapKV & 0.541 & 0.744 & 8\% & 15\% & 8\% & 75\% & 0.181 & --- \\
DropKV & 0.519 & 0.694 & 10\% & 18\% & 8\% & 75\% & 0.199 & --- \\
Ada-KV & 0.731 & 0.879 & 1\% & 15\% & 1\% & 75\% & 0.175 & --- \\
LaProx & 0.755 & 0.915 & 1\% & 11\% & 0\% & 75\% & 0.181 & --- \\
OBCache-K + Ada-KV & 0.774 & 0.914 & 0\% & 10\% & 2\% & 75\% & 0.177 & --- \\
\sieve{} interior only & 0.531 & 0.687 & 34\% & 6\% & 5\% & 65\% & 0.090 & --- \\
\sieve{} (router) & 0.506 & 0.721 & 40\% & 6\% & 1\% & 64\% & 0.089 & 99\% \\
\sieve{} interior, pooled score (diag.) & 0.799 & 0.868 & 1\% & 11\% & 1\% & 64\% & 0.110 & --- \\
\sieve{} oracle router (diag.) & 0.661 & 0.956 & 11\% & 6\% & 1\% & 64\% & 0.085 & 97\% \\
\bottomrule
\end{tabular}

\end{table}

\begin{table}[h]
\centering
\scriptsize
\setlength{\tabcolsep}{2.4pt}
\caption{Paired end-task grid: mean score over 20 prompts per cell, all methods in one process per cell.  Bold marks the best method in each valid row.  $\ddagger$: uncompressed score below 0.9, excluded from every aggregate.  Mistral-7B rows are summarized separately in Table~\ref{tab:endtask}.}
\label{tab:app-grid}
\resizebox{\linewidth}{!}{% generated by figures/make_endtask_tables.py from 9 parquets
\begin{tabular}{@{}llllrrrrrrrrrrrrr@{}}
\toprule
model & ctx & task & $B$ & FP & TurboQ. & KIVI & KIVI-32 & KVQuant & H2O & SnapKV & DropKV & Ada-KV & LaProx & OBC+Ada & \sieve{}-int. & \sieve{} \\
\midrule
Llama-3.1-8B & 8k & multi-key & 2 & 1.00 & \textbf{1.00} & 0.90 & \textbf{1.00} & \textbf{1.00} & 0.20 & 0.15 & 0.25 & 0.30 & 0.30 & 0.30 & 0.35 & 0.40 \\
Llama-3.1-8B & 8k & multi-key & 3 & 1.00 & \textbf{1.00} & \textbf{1.00} & \textbf{1.00} & \textbf{1.00} & 0.30 & 0.40 & 0.30 & 0.55 & 0.70 & 0.55 & 0.65 & 0.95 \\
Llama-3.1-8B & 8k & multi-value & 2 & 1.00 & \textbf{1.00} & 0.90 & \textbf{1.00} & 0.95 & 0.30 & 0.15 & 0.14 & 0.17 & 0.20 & 0.28 & 0.12 & 0.17 \\
Llama-3.1-8B & 8k & multi-value & 3 & 1.00 & \textbf{1.00} & \textbf{1.00} & 0.95 & 0.95 & 0.45 & 0.25 & 0.26 & 0.34 & 0.45 & 0.59 & 0.41 & 0.99 \\
Llama-3.1-8B & 8k & single & 2 & 1.00 & \textbf{1.00} & \textbf{1.00} & \textbf{1.00} & 0.95 & 0.00 & 0.30 & 0.25 & 0.50 & 0.55 & 0.70 & 0.30 & 0.30 \\
Llama-3.1-8B & 8k & single & 3 & 1.00 & \textbf{1.00} & \textbf{1.00} & \textbf{1.00} & \textbf{1.00} & 0.00 & 0.50 & 0.45 & 0.75 & 0.90 & 0.90 & 0.65 & \textbf{1.00} \\
Llama-3.1-8B & 8k & var.\ track. & 2 & 1.00 & \textbf{0.98} & 0.91 & 0.94 & 0.96 & 0.79 & 0.40 & 0.44 & 0.38 & 0.41 & 0.51 & 0.42 & 0.63 \\
Llama-3.1-8B & 8k & var.\ track. & 3 & 1.00 & \textbf{1.00} & \textbf{1.00} & \textbf{1.00} & \textbf{1.00} & 0.79 & 0.63 & 0.56 & 0.72 & 0.81 & 0.82 & 0.83 & \textbf{1.00} \\
\addlinespace[1pt]
Llama-3.1-8B & 32k & multi-key & 2 & 1.00 & \textbf{1.00} & 0.95 & \textbf{1.00} & \textbf{1.00} & 0.20 & 0.15 & 0.20 & 0.40 & 0.45 & 0.45 & 0.45 & 0.55 \\
Llama-3.1-8B & 32k & multi-key & 3 & 1.00 & \textbf{1.00} & \textbf{1.00} & \textbf{1.00} & \textbf{1.00} & 0.25 & 0.55 & 0.25 & 0.65 & 0.75 & 0.65 & 0.65 & \textbf{1.00} \\
Llama-3.1-8B & 32k & multi-value & 2 & 1.00 & 0.96 & 0.97 & \textbf{1.00} & 0.93 & 0.28 & 0.15 & 0.10 & 0.25 & 0.28 & 0.35 & 0.11 & 0.12 \\
Llama-3.1-8B & 32k & multi-value & 3 & 1.00 & \textbf{1.00} & \textbf{1.00} & \textbf{1.00} & \textbf{1.00} & 0.44 & 0.36 & 0.36 & 0.57 & 0.56 & 0.61 & 0.30 & 0.97 \\
Llama-3.1-8B & 32k & single & 2 & 1.00 & \textbf{1.00} & 0.90 & \textbf{1.00} & \textbf{1.00} & 0.05 & 0.50 & 0.40 & 0.75 & 0.75 & 0.80 & 0.35 & 0.45 \\
Llama-3.1-8B & 32k & single & 3 & 1.00 & \textbf{1.00} & \textbf{1.00} & \textbf{1.00} & \textbf{1.00} & 0.05 & 0.75 & 0.65 & 0.95 & 0.80 & 0.95 & 0.60 & \textbf{1.00} \\
Llama-3.1-8B & 32k & var.\ track. & 2 & 1.00 & \textbf{0.98} & 0.90 & 0.97 & \textbf{0.98} & 0.80 & 0.29 & 0.30 & 0.66 & 0.57 & 0.68 & 0.46 & 0.53 \\
Llama-3.1-8B & 32k & var.\ track. & 3 & 1.00 & \textbf{1.00} & \textbf{1.00} & \textbf{1.00} & 0.98 & 0.78 & 0.62 & 0.61 & 0.86 & 0.85 & 0.85 & 0.66 & \textbf{1.00} \\
\addlinespace[1pt]
Llama-3.1-8B & 128k & multi-key & 2 & 1.00 & 0.95 & 0.85 & \textbf{1.00} & 0.95 & 0.10 & 0.45 & 0.35 & 0.55 & 0.65 & 0.65 & 0.45 & 0.40 \\
Llama-3.1-8B & 128k & multi-key & 3 & 1.00 & \textbf{1.00} & \textbf{1.00} & \textbf{1.00} & \textbf{1.00} & 0.25 & 0.60 & 0.65 & 0.80 & 0.90 & 0.90 & 0.55 & 0.55 \\
Llama-3.1-8B & 128k & multi-value & 2 & 0.97 & 0.96 & 0.89 & \textbf{0.97} & 0.89 & 0.23 & 0.47 & 0.42 & 0.68 & 0.70 & 0.74 & 0.51 & 0.46 \\
Llama-3.1-8B & 128k & multi-value & 3 & 0.97 & 0.99 & 0.97 & \textbf{1.00} & 0.96 & 0.34 & 0.64 & 0.59 & 0.82 & 0.85 & 0.82 & 0.64 & 0.72 \\
Llama-3.1-8B & 128k & single & 2 & 1.00 & \textbf{1.00} & \textbf{1.00} & \textbf{1.00} & \textbf{1.00} & 0.05 & 0.80 & 0.80 & 0.95 & 0.95 & 0.95 & 0.50 & 0.45 \\
Llama-3.1-8B & 128k & single & 3 & 1.00 & \textbf{1.00} & \textbf{1.00} & \textbf{1.00} & \textbf{1.00} & 0.10 & 0.95 & 0.80 & \textbf{1.00} & \textbf{1.00} & \textbf{1.00} & 0.70 & 0.70 \\
Llama-3.1-8B & 128k & var.\ track. & 2 & 0.96 & 0.81 & \textbf{0.86} & 0.82 & 0.78 & 0.58 & 0.44 & 0.50 & 0.75 & 0.72 & 0.76 & 0.66 & 0.71 \\
Llama-3.1-8B & 128k & var.\ track. & 3 & 0.96 & 0.95 & 0.96 & \textbf{0.99} & 0.92 & 0.77 & 0.79 & 0.74 & 0.89 & 0.91 & 0.93 & 0.86 & 0.91 \\
\addlinespace[1pt]
Mistral-7B & 8k & multi-key & 2 & 1.00 & 0.95 & 0.95 & 0.95 & \textbf{1.00} & 0.15 & 0.00 & 0.00 & 0.00 & 0.00 & 0.00 & 0.00 & 0.70 \\
Mistral-7B & 8k & multi-key & 3 & 1.00 & \textbf{1.00} & 0.95 & \textbf{1.00} & \textbf{1.00} & 0.25 & 0.00 & 0.15 & 0.05 & 0.10 & 0.05 & 0.15 & \textbf{1.00} \\
Mistral-7B & 8k & multi-value & 2 & 0.93 & \textbf{0.76} & 0.17 & 0.55 & 0.28 & 0.11 & 0.00 & 0.00 & 0.00 & 0.00 & 0.00 & 0.00 & 0.45 \\
Mistral-7B & 8k & multi-value & 3 & 0.93 & 0.90 & 0.62 & \textbf{0.95} & 0.49 & 0.35 & 0.00 & 0.01 & 0.01 & 0.01 & 0.04 & 0.00 & 0.94 \\
Mistral-7B & 8k & single & 2 & 0.95 & \textbf{1.00} & \textbf{1.00} & 0.95 & \textbf{1.00} & 0.00 & 0.00 & 0.00 & 0.10 & 0.10 & 0.15 & 0.05 & 0.60 \\
Mistral-7B & 8k & single & 3 & 0.95 & \textbf{1.00} & 0.95 & 0.95 & 0.95 & 0.00 & 0.10 & 0.25 & 0.30 & 0.35 & 0.25 & 0.30 & 0.95 \\
Mistral-7B & 8k & var.\ track.$^\ddagger$ & 2 & 0.20 & 0.20 & 0.24 & 0.20 & 0.20 & 0.20 & 0.04 & 0.06 & 0.06 & 0.10 & 0.13 & 0.09 & 0.19 \\
Mistral-7B & 8k & var.\ track.$^\ddagger$ & 3 & 0.20 & 0.20 & 0.20 & 0.20 & 0.20 & 0.20 & 0.10 & 0.16 & 0.15 & 0.15 & 0.18 & 0.15 & 0.20 \\
\addlinespace[1pt]
Mistral-7B & 32k & multi-key & 2 & 1.00 & 0.95 & 0.85 & \textbf{1.00} & 0.90 & 0.05 & 0.00 & 0.05 & 0.00 & 0.00 & 0.05 & 0.00 & 0.10 \\
Mistral-7B & 32k & multi-key & 3 & 1.00 & \textbf{1.00} & \textbf{1.00} & \textbf{1.00} & \textbf{1.00} & 0.05 & 0.00 & 0.05 & 0.05 & 0.05 & 0.10 & 0.05 & 0.95 \\
Mistral-7B & 32k & multi-value$^\ddagger$ & 2 & 0.14 & 0.07 & 0.04 & 0.23 & 0.04 & 0.00 & 0.00 & 0.00 & 0.00 & 0.00 & 0.01 & 0.00 & 0.04 \\
Mistral-7B & 32k & multi-value$^\ddagger$ & 3 & 0.14 & 0.04 & 0.04 & 0.14 & 0.04 & 0.05 & 0.03 & 0.06 & 0.00 & 0.00 & 0.00 & 0.00 & 0.04 \\
Mistral-7B & 32k & single & 2 & 1.00 & 0.95 & 0.75 & 0.95 & \textbf{1.00} & 0.00 & 0.05 & 0.15 & 0.10 & 0.25 & 0.25 & 0.00 & 0.30 \\
Mistral-7B & 32k & single & 3 & 1.00 & \textbf{1.00} & \textbf{1.00} & \textbf{1.00} & \textbf{1.00} & 0.00 & 0.10 & 0.20 & 0.40 & 0.45 & 0.50 & 0.15 & \textbf{1.00} \\
Mistral-7B & 32k & var.\ track.$^\ddagger$ & 2 & 0.21 & 0.20 & 0.20 & 0.20 & 0.20 & 0.18 & 0.07 & 0.07 & 0.11 & 0.12 & 0.14 & 0.11 & 0.18 \\
Mistral-7B & 32k & var.\ track.$^\ddagger$ & 3 & 0.21 & 0.20 & 0.20 & 0.20 & 0.20 & 0.22 & 0.18 & 0.14 & 0.16 & 0.16 & 0.17 & 0.14 & 0.20 \\
\addlinespace[1pt]
Qwen3-8B & 8k & multi-key & 2 & 1.00 & 0.95 & \textbf{1.00} & \textbf{1.00} & \textbf{1.00} & 0.25 & 0.05 & 0.00 & 0.00 & 0.10 & 0.05 & 0.40 & 0.90 \\
Qwen3-8B & 8k & multi-key & 3 & 1.00 & \textbf{1.00} & \textbf{1.00} & \textbf{1.00} & \textbf{1.00} & 0.35 & 0.10 & 0.20 & 0.15 & 0.55 & 0.25 & 0.65 & \textbf{1.00} \\
Qwen3-8B & 8k & multi-value$^\ddagger$ & 2 & 0.25 & 0.79 & 0.23 & 0.39 & 0.33 & 0.10 & 0.01 & 0.04 & 0.00 & 0.10 & 0.03 & 0.06 & 0.80 \\
Qwen3-8B & 8k & multi-value$^\ddagger$ & 3 & 0.25 & 0.40 & 0.49 & 0.20 & 0.46 & 0.12 & 0.04 & 0.06 & 0.04 & 0.30 & 0.06 & 0.23 & 0.49 \\
Qwen3-8B & 8k & single & 2 & 1.00 & \textbf{1.00} & 0.95 & \textbf{1.00} & \textbf{1.00} & 0.00 & 0.10 & 0.20 & 0.10 & 0.25 & 0.10 & 0.30 & 0.95 \\
Qwen3-8B & 8k & single & 3 & 1.00 & \textbf{1.00} & \textbf{1.00} & \textbf{1.00} & \textbf{1.00} & 0.00 & 0.25 & 0.35 & 0.25 & 0.70 & 0.30 & 0.45 & \textbf{1.00} \\
Qwen3-8B & 8k & var.\ track. & 2 & 1.00 & 0.42 & 0.98 & \textbf{1.00} & 0.94 & 0.74 & 0.17 & 0.29 & 0.26 & 0.40 & 0.40 & 0.43 & 0.87 \\
Qwen3-8B & 8k & var.\ track. & 3 & 1.00 & \textbf{1.00} & 0.99 & \textbf{1.00} & \textbf{1.00} & 0.81 & 0.55 & 0.53 & 0.64 & 0.73 & 0.69 & 0.82 & \textbf{1.00} \\
\addlinespace[1pt]
Qwen3-8B & 32k & multi-key & 2 & 1.00 & 0.70 & 0.95 & \textbf{1.00} & \textbf{1.00} & 0.10 & 0.10 & 0.15 & 0.15 & 0.25 & 0.15 & 0.20 & 0.90 \\
Qwen3-8B & 32k & multi-key & 3 & 1.00 & \textbf{1.00} & \textbf{1.00} & \textbf{1.00} & \textbf{1.00} & 0.25 & 0.20 & 0.35 & 0.20 & 0.45 & 0.20 & 0.50 & 0.95 \\
Qwen3-8B & 32k & multi-value$^\ddagger$ & 2 & 0.00 & 0.82 & 0.25 & 0.14 & 0.21 & 0.03 & 0.00 & 0.05 & 0.00 & 0.00 & 0.04 & 0.00 & 0.19 \\
Qwen3-8B & 32k & multi-value$^\ddagger$ & 3 & 0.00 & 0.24 & 0.04 & 0.04 & 0.09 & 0.03 & 0.01 & 0.00 & 0.01 & 0.00 & 0.00 & 0.05 & 0.04 \\
Qwen3-8B & 32k & single & 2 & 1.00 & \textbf{1.00} & \textbf{1.00} & \textbf{1.00} & \textbf{1.00} & 0.00 & 0.15 & 0.30 & 0.15 & 0.20 & 0.20 & 0.25 & \textbf{1.00} \\
Qwen3-8B & 32k & single & 3 & 1.00 & \textbf{1.00} & \textbf{1.00} & \textbf{1.00} & \textbf{1.00} & 0.00 & 0.25 & 0.50 & 0.20 & 0.50 & 0.30 & 0.55 & \textbf{1.00} \\
Qwen3-8B & 32k & var.\ track. & 2 & 1.00 & 0.44 & 0.96 & \textbf{0.98} & 0.90 & 0.77 & 0.20 & 0.36 & 0.33 & 0.53 & 0.53 & 0.37 & \textbf{0.98} \\
Qwen3-8B & 32k & var.\ track. & 3 & 1.00 & 0.98 & \textbf{1.00} & \textbf{1.00} & \textbf{1.00} & 0.81 & 0.54 & 0.69 & 0.68 & 0.78 & 0.70 & 0.64 & 0.99 \\
\bottomrule
\end{tabular}
}
\end{table}



\section{Extended related work}
\label{app:related}

\begin{table}[h]
\centering
\footnotesize
\setlength{\tabcolsep}{3pt}
\caption{Closest work and the claim boundary of this paper.}
\label{tab:related}
\begin{tabular}{@{}p{0.19\linewidth}p{0.29\linewidth}p{0.45\linewidth}@{}}
\toprule
\textbf{Work} & \textbf{Allocation/action} & \textbf{Difference here} \\
\midrule
DiffKV; HqeKV & token high/low precision and pruning & explain and predict when the middle tier pays; no first-hybrid claim \\
CAOTE; CriticalKV; DropKV; OBCache; LaProx; ReST-KV & output-aware token eviction & finite-rate extension under a stated local surrogate; exact set error retained; evaluated as baselines where available \\
RateQuant & K/V heads, finite rates & within-head key tokens, explicit zero rate, context traversal \\
RDKV  & value tokens, key channels, zero-to-finite rates & when the token-wise interior beats endpoints; score staleness and regime movement \\
AATC  & transformed K/V channels, including zero & token-axis tier viability; values/channels remain outside our allocation \\
Fixed-contract diagnostic & attribution of eviction failures & measures how comparison set and information contract move verdicts; truncation mechanism \\
\bottomrule
\end{tabular}
\end{table}

\paragraph{Low-precision formats and the dense endpoint.}
Block-scaled 4-bit formats strengthen the dense endpoint.  NVFP4 pairs 4-bit
E2M1 values with FP8 scales over 16-element blocks~\citep{alvarez2025nvfp4,nvidia2025nvfp4ways},
finer than the power-of-two, 32-element scales of MXFP4~\citep{ocp2023mx}, and
libraries provide fused FP8/FP4 kernels~\citep{nvidia2024te,nvidia2024modelopt}.
A 4-bit cache with hardware support is therefore the baseline any mixed or
evicting policy must beat.
% Commented out (2026-09-26): broader FP4 context and the untested per-channel
% outlier hypothesis for TurboQuant on Qwen3.  Original paragraph:
% Block-scaled FP4 has moved from storage compression to full-stack inference and
% training.  NVFP4 pairs 4-bit E2M1 values with FP8 scales over 16-element
% blocks~\citep{alvarez2025nvfp4,nvidia2025nvfp4ways}, finer than the
% power-of-two, 32-element scales of MXFP4~\citep{ocp2023mx}, and has been used
% to pretrain LLMs~\citep{nvidia2025nvfp4pretrain,devleker2025nvfp4train}.
% Libraries provide fused FP8/FP4 kernels and quantization
% recipes~\citep{nvidia2024te,nvidia2024modelopt}.  Open-weight models ship with
% FP4 expert weights~\citep{openai2025gptoss,deepseek2026v4}, FP4 has been
% applied to diffusion models by absorbing outliers into low-rank
% branches~\citep{li2025svdquant} and to diffusion transformers for image and
% video~\citep{zhao2024viditq}, and NVFP4 backbones have been combined with LoRA
% for reinforcement learning~\citep{huang2025qerl}.  For our question, these
% formats strengthen the dense endpoint: a 4-bit cache with hardware support is
% the baseline any mixed or evicting policy must beat.  Its 16-element scales
% also limit how far a few large key channels can inflate the quantization step,
% the weakness of per-token scaling that we suspect behind TurboQuant's failure
% on Qwen3 (F4); we have not tested NVFP4 caches here.

\paragraph{Low-precision attention kernels.}
FlashAttention-2 and -3 fuse attention without materializing the score
matrix, the latter with FP8~\citep{dao2023flashattention2,shah2024flashattention3};
SageAttention quantizes the attention computation itself to 8 bits, INT4 with
outlier smoothing, and microscaled FP4~\citep{zhang2024sageattention,zhang2024sageattention2,zhang2025sageattention3}.
Because these kernels never write attention scores to memory, evictors that
read historical scores (H2O's accumulated attention, and the lagged scores our
allocator uses) must either fall back to eager attention or estimate
importance another way, for example from pre-RoPE geometry~\citep{mao2026triattention}.
Paged allocators also reclaim only empty blocks~\citep{kwon2023vllm}, so
scattered surviving tokens free little memory unless the cache is
compacted~\citep{mao2026triattention}.
Our 4-bit cascade avoids this: it re-scores tokens from stored 4-bit keys in a
small separate pass over the 32 window queries, so it needs no scores from the
fused kernel.
Our experiments use an eager attention path and therefore measure accuracy,
not these costs.
